\documentclass[lettersize,journal]{IEEEtran}
\usepackage{amsmath,amsfonts}
\usepackage{algorithmic}
\usepackage{algorithm}
\usepackage{array}
\usepackage[caption=false,font=normalsize,labelfont=sf,textfont=sf]{subfig}
\usepackage{textcomp}
\usepackage{stfloats}
\usepackage{url}
\usepackage{verbatim}
\usepackage{graphicx}
\usepackage{cite}
\hyphenation{op-tical net-works semi-conduc-tor IEEE-Xplore}

\begin{document}

\title{High-Performance Telemetry Visualization: Native C++ Architectures vs. Standard Web-Stack Baselines in Factory Workstations}

\author{Divyam Agarwal%
\thanks{Divyam Agarwal is with Delhi Technological University (e-mail: iitdivyamagarwal@gmail.com).}}

% The paper headers
\markboth{High-Performance Telemetry Visualization}%
{Agarwal: High-Performance Telemetry Visualization}

\maketitle

\begin{abstract}
In modern Industry 4.0 environments, factory monitoring workstations (``thick clients'') are tasked with visualizing high-frequency, asynchronous telemetry data in real-time. Despite having adequate computing and graphics resources, developers frequently default to deploying Web-based dashboards wrapped in browser engines. While a naive DOM/WebSocket implementation introduces severe frame-rate instability, we demonstrate that even the best-optimized Full-Stack Web client (Node.js + Canvas) cannot match a Native C++ architecture on any performance axis. By expanding our evaluation to include complete process-tree CPU measurements and precise Proportional Set Size (PSS) memory accounting, we provide a rigorous comparison of these architectures under a 5,000 MQTT messages/second workload. Our results show that while all architectures successfully delivered 100\% of the published payload to the subscriber application, the Native C++ (GTK4) architecture achieved the lowest CPU utilization ($27.4 \pm 1.2\%$), the smallest memory footprint ($57.4$ MB PSS), and the lowest end-to-end median display latency ($0.2 \pm 0.0$ ms). In contrast, the Full-Stack Web architecture required 44\% more CPU and 59\% more memory while delivering 4.5$\times$ higher median latency ($0.9 \pm 0.1$ ms). Native GTK4 reclaims 34 MB of PSS per interface, making it the strict preference for environments with constrained RAM headroom, such as those running concurrent edge-AI workloads.
\end{abstract}

\begin{IEEEkeywords}
Edge Computing, Factory Automation, GTK4, Industrial Internet of Things, MQTT, Node.js, Python, Telemetry, WebSockets.
\end{IEEEkeywords}

\section{Introduction}
\label{sec:introduction}
\IEEEPARstart{T}{he} monitoring of real-time factory floor data requires interfaces that are highly responsive and capable of absorbing massive data floods without lagging or crashing. As Industry 4.0 relies increasingly on edge computing for real-time industrial automation \cite{b1, b4}, telemetry systems have had to evolve to handle sustained, high-frequency continuous data streams.

A common trend in industry is to build these interfaces using rapid-development web frameworks (like Python Flask or Node.js) and display them locally on a workstation via a web browser. While easy to develop, a naive DOM-based implementation of this ``web-stack on a workstation'' introduces severe rendering latency. As our results will show, the performance of web-based frameworks can be significantly improved by switching to faster asynchronous runtimes (e.g., Node.js) and HTML5 Canvas, but they still incur a systemic memory and CPU overhead due to the fixed architectural cost of the multi-process browser engine \cite{b5}.

Furthermore, modern factory monitoring workstations are frequently required to run heavy concurrent workloads, such as real-time object detection models (e.g., YOLOv8 \cite{b2} when run in CPU-only edge deployments) for automated industrial quality inspection \cite{b3}. Because these AI quality-control workloads demand maximum computational resources, it is imperative that the telemetry visualization interface consumes as little overhead as possible. This paper evaluates whether a return to highly optimized, Native C++ architectures (utilizing GTK4) can bypass the browser memory tax and reclaim critical hardware headroom, and compares the resulting performance against state-of-the-art web architectures.

\subsection{Reproducibility and Open Data}
To ensure transparency and reproducibility, all benchmark scripts, the MQTT publisher, the C++ application source, the Node.js implementation, and raw measurement logs for the $N=10$ trials are available at: \url{https://github.com/Inuth0603/flask-mqtt-vs-native}.

\section{Methodology}

To evaluate the architectural efficiency of telemetry visualization, we simulated a high-frequency data flood representing factory sensor output. We compared a Native C++ Interface against three Web Interfaces: a standard DOM implementation (Flask backend), a Canvas implementation (Flask backend), and a Full-Stack Optimized implementation (Node.js backend + Canvas).

\subsection{Experimental Setup \& Effective Workload Verification}
To ensure strict statistical validity, we conducted $N=10$ independent trials for each architecture. Each trial consisted of a 60-second sustained data flood. A single Python publisher asynchronously published packets to an Eclipse Mosquitto broker at a pinned rate of \textbf{5,000 MQTT messages per second} (QoS 0, variable-length plain-text payloads of $\sim$40-50 bytes), resulting in 300,000 total published messages per trial. The publisher, broker, and subscriber all ran concurrently on the same host. To ensure all architectures processed the same effective workload, we verified end-to-end delivery at the subscriber layer. Across all $N=10$ trials, the backend subscribers successfully received exactly 300,000 messages (a 100\% delivery rate), confirming that differences in resource consumption are strictly attributable to application architecture rather than unequal data ingestion.

All benchmarks were executed on a standardized testbed: a 12th Gen Intel Core i5-12450H processor (8 physical cores / 12 siblings), 14 GiB accessible RAM, running Ubuntu natively on a Wayland display server. The Wayland environment utilized the default GSK OpenGL renderer for GTK4. Browser benchmarks were executed on Chromium. Proportional Set Size (PSS) was calculated by aggregating the \texttt{smaps\_rollup} files for all PIDs within the target architecture's process tree. The testbed utilized a 144Hz display, which the Chromium engine fully saturated, whereas the Native GTK4 application targeted a 60Hz update rate via an intentional decoupled polling design.

\subsection{Absolute Wall-Clock Timestamping \& Latency}
To calculate end-to-end latency, the pipeline relies strictly on \textbf{absolute wall-clock Epoch time}, avoiding monotonic timers whose zero-points are arbitrary and differ across disparate runtimes. Because the publisher, broker, and subscriber were executed on the same host machine, all timestamps reference a single shared system clock, eliminating network time synchronization errors.
\begin{itemize}
    \item \textbf{The Publisher:} End-to-end latency is measured from the absolute epoch timestamp applied by the Python publisher (\texttt{time.time()} with microsecond resolution) prior to broker transmission, to the execution of the frontend render callback.
    \item \textbf{The Web Frontend (Browser):} The browser calculates the display timestamp using the W3C High-Resolution Time API: \texttt{performance.timeOrigin + performance.now()}. This yields sub-millisecond precision tied securely to the absolute epoch, captured at the start of \texttt{requestAnimationFrame}.
    \item \textbf{The Native Frontend (GTK4):} The C++ application captures the display timestamp using GLib's \texttt{g\_get\_real\_time() / 1000.0} at \texttt{queue\_draw}, which then waits for the Wayland compositor to pick it up.
\end{itemize}

\subsection{Web Baseline Configurations}
The Flask baseline was deployed using an Eventlet asynchronous worker pool. Standard \texttt{eventlet.monkey\_patch()} was applied at application startup to ensure the Python standard library supported non-blocking I/O for WebSockets. The Node.js baseline utilized the native \texttt{ws} library. Both optimized web variants utilized HTML5 Canvas rendering.

\subsection{Accurate Resource Accounting}
To address common pitfalls in benchmarking multi-process browser applications, our methodology incorporates two critical refinements:
1) \textbf{Comprehensive CPU Measurement:} CPU metrics were sampled at 1-second intervals and normalized to a single core (0--100\%). The measured process tree encompasses only the subscriber/rendering application (e.g., the C++ binary or the Node process plus all Chromium helper threads); the publisher and Mosquitto broker were excluded to isolate visualization overhead.
2) \textbf{Proportional Memory Accounting:} Instead of relying on Resident Set Size (RSS), which can count shared libraries and memory pages multiple times across browser sub-processes \cite{b5}, we utilize Proportional Set Size (PSS) and Unique Set Size (USS) to accurately quantify the true physical memory consumption of each architecture.

\section{Native Architecture \& Thread Decoupling}

GTK4 on Wayland defaults to the monitor's refresh rate (144Hz). The 61.4 FPS measured in our Native architecture is the result of an intentional decoupled polling design, meaning the native app coalesces data while the web apps process every single frame.

\subsection{The GTK4 Update Loop}
\begin{itemize}
    \item \textbf{The Native Pipeline (C++):} A background network thread receives the 5,000 MQTT messages/sec. Upon receiving a payload, it acquires a \texttt{std::mutex}, overwrites a shared \texttt{TelemetryState} struct with the latest values (effectively dropping/coalescing intermediate messages), and releases the lock.
    \item \textbf{The Draw Trigger:} The main GTK UI thread runs a GLib timer loop (\texttt{g\_timeout\_add}) set to a 16ms interval (targeting 60 FPS). Every 16ms, it wakes up, locks the mutex, reads the latest state, and calls \texttt{gtk\_widget\_queue\_draw()}.
\end{itemize}

\subsection{Latency Characteristics}
With a correctly functioning network thread (see Section \ref{sec:introduction}, Reproducibility), the 16 ms GLib polling timer yields sub-millisecond display latency (0.2 ms median, 0.4 ms P99). At 5,000 messages/sec, a new payload arrives on average every 0.2 ms. Since the UI thread polls at 16 ms intervals, it virtually always finds fresh data in the shared buffer. The display latency therefore reflects only the time from the most recent \texttt{message\_arrived()} write to the next \texttt{CLOCK\_MONOTONIC} read in the timer callback—dominated by mutex acquisition and a single memory copy, both of which complete in sub-microsecond time. The tight standard deviation ($\pm 0.0$ ms) across all 10 trials confirms highly deterministic scheduling under the Wayland compositor.

\section{Results}

The Native C++ architecture demonstrated clear superiority across all measured performance axes. Averages recorded across $N=10$ trials are presented in Table \ref{tab_metrics}.

\begin{table*}[t]
\caption{Comparative Performance Metrics (Mean $\pm$ SD)}
\label{tab_metrics}
\begin{center}
\begin{tabular}{|p{130pt}|p{75pt}|p{85pt}|p{85pt}|p{85pt}|}
\hline
\textbf{Metric} & \textbf{Native C++ (GTK4)} & \textbf{Naive Web (Flask/DOM)*} & \textbf{Optimized Web (Flask/Canvas)} & \textbf{Full-Stack Web (Node.js/Canvas)} \\
\hline
\textbf{Total Process-Tree CPU (\%)} & $27.4 \pm 1.2$ & $33.5 \pm 0.8$* & $35.9 \pm 1.1$ & $39.4 \pm 1.2$ \\
\hline
\textbf{Total Memory (PSS)} & $57.4 \text{ MB} \pm 0.0 \text{ MB}$ & $105.9 \text{ MB} \pm 0.1 \text{ MB}$* & $105.6 \text{ MB} \pm 0.2 \text{ MB}$ & $91.0 \text{ MB} \pm 6.1 \text{ MB}$ \\
\hline
\textbf{Total Memory (USS)} & $35.9 \text{ MB} \pm 0.0 \text{ MB}$ & $102.3 \text{ MB} \pm 0.1 \text{ MB}$* & $101.9 \text{ MB} \pm 0.2 \text{ MB}$ & $88.0 \text{ MB} \pm 6.1 \text{ MB}$ \\
\hline
\textbf{Median End-to-End Latency} & $0.2 \text{ ms} \pm 0.0 \text{ ms}$ & N/A (Lag) & $38.4 \text{ ms} \pm 14.2 \text{ ms}$ & $0.9 \text{ ms} \pm 0.1 \text{ ms}$ \\
\hline
\textbf{99th-Percentile (P99) Latency} & $0.4 \text{ ms} \pm 0.0 \text{ ms}$ & N/A (Lag) & $44.5 \text{ ms} \pm 2.1 \text{ ms}$ & $1.7 \text{ ms} \pm 0.1 \text{ ms}$ \\
\hline
\textbf{Measured Frame Rate (FPS)} & $61.4 \text{ FPS} \pm 0.1$ & $<$10 FPS & $144.0 \text{ FPS} \pm 0.0$ & $144.0 \text{ FPS} \pm 0.0$ \\
\hline
\multicolumn{5}{p{450pt}}{*Metrics captured during system failure (UI thread starvation); provided for baseline context only.}
\end{tabular}
\end{center}
\end{table*}

\subsection{Memory Efficiency (PSS and USS)}
By utilizing Proportional Set Size (PSS), we effectively account for Chromium's shared memory pages. The Native C++ application demonstrates superior memory efficiency, consuming only 57.4 MB (PSS) and 35.9 MB (USS) under peak load. In contrast, the Full-Stack Node.js web architecture required 91.0 MB (PSS), representing a 59\% increase in true memory footprint due to the multi-process browser architecture. Welch's t-test confirmed a statistically significant difference in memory footprint between the Native and Node.js implementations ($t(9.0) = -17.37, p < 0.001$). The Flask variants exhibited even higher memory consumption at approximately 106 MB PSS, owing to the additional overhead of the Python/Eventlet runtime. For resource-constrained edge devices, the native approach significantly preserves available RAM. Industrial edge devices like the Siemens IPC127E typically operate with only $\sim$2.2 GB of free RAM for user applications \cite{b6}. In these highly constrained environments, saving 34 MB of PSS per telemetry interface is operationally critical, and the advantage compounds with each additional dashboard instance.

\subsection{CPU Utilization and End-to-End Latency}
The Native C++ implementation achieved the lowest CPU utilization across all architectures at $27.4 \pm 1.2\%$ of a single core, while simultaneously delivering the lowest display latency (0.2 ms median, 0.4 ms P99). This result underscores the efficiency of the direct MQTT-to-GTK4 rendering pipeline: the background network thread writes data into shared memory with minimal contention, and the 60Hz polling timer ensures the coalesced state is dispatched to the Wayland compositor with sub-millisecond overhead.

The Node.js Canvas implementation also achieved sub-millisecond responsiveness, posting a median end-to-end latency of 0.9 ms (and 1.7 ms P99) while sustaining 144 FPS rendering. However, this came at a 44\% higher CPU cost ($39.4 \pm 1.2\%$), attributable to the compounded overhead of the multi-process Chromium engine and the unrestrained 144 FPS rendering rate. Even the Flask variants, despite their Python/Eventlet runtime, consumed more CPU ($33.5$--$35.9\%$) than the native binary, while the Flask Canvas variant exhibited significantly degraded latency (38.4 ms median) due to the additional WebSocket hop through the Python backend.

\subsection{The Cost of String Concatenation (Flask DOM)}
The Naive Web implementation (Flask DOM) experienced severe lag, dropping below 10 FPS and failing to record accurate latency metrics. This failure exposes a classic frontend performance bottleneck: continuously mutating the browser DOM (e.g., rewriting an unbounded text element) at 5,000 updates per second incurs massive reflow and string concatenation overhead, causing the main rendering thread to saturate. Moving to Canvas rendering fully mitigates this issue, as seen in the Flask Canvas variant.

\section{Threats to Validity}
While our results provide clear distinctions between architectures, the evaluation is based on a single fixed workload (5,000 msg/sec, QoS 0, 60s duration). Architectural behavior may differ under varying payload sizes, increased burst patterns, or different QoS settings. Furthermore, benchmarks were executed on a native Wayland environment; different display servers (e.g., X11) or embedded runtimes with differing GPU-acceleration profiles could alter the CPU/Latency balance between GTK4 and Chromium rendering. Future work should systematically evaluate these architectures across varying network conditions and hardware profiles.

\section{Conclusion}
This paper demonstrates a clear architectural advantage for Native C++ telemetry visualization under sustained Industry 4.0 workloads. When subjected to a 5,000 msg/sec data flood, the Native GTK4 architecture dominated across all measured performance axes: it consumed the least CPU ($27.4\%$ vs.\ $39.4\%$ for Node.js), maintained the smallest memory footprint ($57.4$ MB PSS vs.\ $91.0$ MB), and achieved the lowest end-to-end display latency ($0.2$ ms median vs.\ $0.9$ ms). These advantages stem from bypassing the multi-process browser engine entirely, eliminating the systemic overhead of Chromium's renderer, GPU, and utility processes. Native GTK4 reclaims 34 MB of PSS per interface, making it the strict preference for environments with constrained RAM headroom, particularly edge workstations running concurrent AI inference workloads.

\section*{Acknowledgments}
This work was supported by independent research and the open-source community. The author would like to thank the developers of the \texttt{Flask-MQTT} library and the broader GTK and Node.js ecosystems for their invaluable tools and documentation.

\begin{thebibliography}{6}
\bibliographystyle{IEEEtran}

\bibitem{b1} K. R. Priya \textit{et al.}, ``Edge Computing Powered Embedded Vision System for Real-Time Industrial Automation and Fault Detection,'' in \textit{2025 International Conference on Metaverse and Current Trends in Computing (ICMCTC)}, IEEE, 2025.

\bibitem{b2} G. Jocher, A. Chaurasia, and J. Qiu, \textit{Ultralytics YOLOv8} (Version 8.0.0) [Computer software], 2023. Available: \url{https://github.com/ultralytics/ultralytics}

\bibitem{b3} T. Gao \textit{et al.}, ``Optimization of Industrial Quality Inspection Systems in Computer Vision: Integration of Deep Learning and IoT Technologies,'' \textit{Journal of Organizational and End User Computing}, vol. 37, IGI Global, 2025.

\bibitem{b4} A. Al-Fuqaha, M. Guizani, M. Mohammadi, M. Aledhari, and M. Ayyash, ``Internet of Things: A Survey on Enabling Technologies, Protocols, and Applications,'' \textit{IEEE Communications Surveys \& Tutorials}, vol. 17, no. 4, pp. 2347--2376, 2015.

\bibitem{b5} C. Reis and S. D. Gribble, ``Isolating Web Programs in Modern Browser Architectures,'' in \textit{Proceedings of the 4th ACM European Conference on Computer Systems (EuroSys)}, 2009, pp. 219--232.

\bibitem{b6} S. S. M., A. Bhadauria, K. Nandy, and S. Upadhyay, ``Lightweight Data Storage and Caching Solution for MQTT Broker on Edge - A Case Study with SQLite and Redis,'' in \textit{2024 IEEE 21st International Conference on Software Architecture Companion (ICSA-C)}, IEEE, 2024, pp. 368--372.

\end{thebibliography}

\vspace{11pt}

\begin{IEEEbiographynophoto}{Divyam Agarwal}
is a second-year undergraduate student in the Department of Production and Industrial Engineering at Delhi Technological University. His current research interests include high-performance computing, systems architecture, and real-time industrial automation.
\end{IEEEbiographynophoto}

\vfill

\end{document}
